\title{The Era of 1-bit LLMs: All Large Language Models are in 1.58 Bits}

\begin{abstract}
Recent advances, exemplified by BitNet~\cite{bitnet}, are ushering in a novel generation of 1-bit Large Language Models (LLMs). In this study, we present a 1-bit LLM variant, called \textbf{\text{BitNet b1.58}}, where every individual parameter (or weight) in the LLM is ternary-valued \{-1, 0, 1\}. This model achieves parity with the full-precision (i.e., FP16 or BF16) Transformer LLM of identical size and trained on the same amount of tokens, both in terms of perplexity and downstream task performance, while offering substantial improvements in latency, memory usage, throughput, and energy efficiency. More importantly, the 1.58-bit LLM introduces a new scaling law and training methodology for developing next-generation LLMs that are both highly effective and resource-efficient. Additionally, it supports a novel computational paradigm and paves the way for designing specialized hardware tailored for 1-bit LLMs.
\end{abstract}

\section{The Era of 1-bit LLMs}

Over recent years, AI has witnessed rapid expansion in the scale and capability of Large Language Models (LLMs). These models have achieved outstanding results across diverse natural language processing tasks, yet their growing size has introduced deployment challenges and sparked concerns regarding their environmental and financial costs due to elevated energy consumption.

One strategy to mitigate these issues is post-training quantization, which creates low-bit models for inference~\cite{smoothquant, gptq, quip, quip_sharp}. This approach lowers the precision of weights and activations, greatly diminishing the memory and computational demands of LLMs. The trend has progressed from 16-bit representations down to lower bit-widths, including 4-bit versions~\cite{gptq, awq}. Nevertheless, despite its widespread adoption in industrial LLM deployments, post-training quantization remains a sub-optimal solution.

Recent advancements in 1-bit model architectures, such as BitNet~\cite{bitnet}, offer a promising avenue for lowering the expense of large language models (LLMs) while preserving their effectiveness. Standard LLMs typically operate with 16-bit floating-point values (i.e., FP16 or BF16), and the primary operation in these models is matrix multiplication. Consequently, the main computational expense arises from floating-point addition and multiplication operations. In contrast, BitNet’s matrix multiplication uses only integer addition, which drastically reduces the energy consumption for LLMs. Since power is often the key limiting factor for computational performance on many chips, these energy savings can translate into faster processing.

Besides computation, transferring model parameters from DRAM to the memory of an on-chip accelerator (such as SRAM) can be costly during inference. While efforts have been made to increase SRAM capacity to boost throughput, this results in significantly higher expenses compared to DRAM. Relative to full-precision models, 1-bit LLMs demand much less memory in terms of both capacity and bandwidth, which can substantially cut down the cost and duration of loading weights from DRAM, leading to quicker and more efficient inference.

In this paper, we propose an important 1-bit LLM variant named \textbf{\text{BitNet b1.58}}, in which each parameter is ternary, assuming values in \{-1, 0, 1\}. By incorporating an additional zero value into the original 1-bit BitNet scheme, we achieve 1.58 bits in binary representation. \text{BitNet b1.58} preserves all advantages of the initial 1-bit BitNet, including its novel computation framework that almost entirely eliminates multiplication in matrix multiplication and allows for significant optimization. Moreover, it maintains the same energy consumption as the original 1-bit BitNet and demonstrates substantially improved efficiency in memory usage, throughput, and latency compared to FP16 LLM baselines. Additionally, \text{BitNet b1.58} provides two further benefits. First, its modeling power is enhanced through explicit feature filtering, enabled by the introduction of the zero value in the model weights, which can greatly boost the performance of 1-bit LLMs. Second, our experimental results indicate that \text{BitNet b1.58} can match full-precision (i.e., FP16) baselines in perplexity and downstream task performance beginning at a 3B parameter scale, using identical configurations (e.g., model size, training tokens, etc.).

\section{BitNet b1.58}

\text{BitNet b1.58} builds upon the BitNet framework, which is a Transformer variant that substitutes the \emph{nn.Linear} layers with \emph{BitLinear}. It is trained from the ground up using weights quantized to 1.58 bits and activations quantized to 8 bits. In comparison to the original BitNet, it incorporates several changes that we outline below.

\paragraph{Quantization Function.} To restrict the weights to values in \{-1, 0, +1\}, we employ an \emph{absmean} quantization strategy. This method first normalizes the weight matrix by its mean absolute value, and then rounds each element to the nearest integer among \{-1, 0, +1\}:

\begin{equation}
    \widetilde{W} = \mathrm{RoundClip}\left(\frac{W}{\gamma + \epsilon}, -1, 1\right),
\end{equation}

\begin{equation}
    \mathrm{RoundClip}(x, a, b) = \max(a, \min(b, \mathrm{round}(x))),
\end{equation}

\begin{equation}
    \gamma = \frac{1}{nm}\sum_{ij} |W_{ij}|.
\end{equation}

For activations, the quantization approach follows that of BitNet except that we do not rescale activations before applying nonlinearities to the $[0, Q_b]$ interval. Instead, activations are scaled per token to the symmetric range $[-Q_b, Q_b]$ to eliminate the need for zero-point quantization. This simplification benefits implementation and system-level optimizations, while causing negligible impact on performance based on our experiments.

\paragraph{LLaMA-alike Components.}

The LLaMA~\cite{llama, llama2} architecture has become the standard foundation for open-source large language models. To support the open-source ecosystem, our design of \text{BitNet b1.58} incorporates LLaMA-inspired components. Concretely, it employs RMSNorm~\cite{rmsnorm}, SwiGLU~\cite{swiglu}, rotary embeddings~\cite{rope}, and removes all bias terms. This allows \text{BitNet b1.58} to be easily integrated into widely used open-source frameworks (e.g., Huggingface, vLLM~\cite{vllm}, and llama.cpp\footnote{\url{https://github.com/ggerganov/llama.cpp}}) with minimal modification.

\section{Results}

We conducted a comparison between \text{BitNet b1.58} and our reproduced FP16 \text{LLaMA LLM} across different model sizes. To maintain fairness, both models were pre-trained on the RedPajama dataset~\cite{redpajama} for 100 billion tokens. The zero-shot capabilities were assessed on several language tasks, such as ARC-Easy~\cite{arc}, ARC-Challenge~\cite{arc}, Hellaswag~\cite{hellaswag}, Winogrande~\cite{winoGrande}, PIQA~\cite{piqa}, OpenbookQA~\cite{openbookqa}, and BoolQ~\cite{boolq}. Additionally, we measured the validation perplexity on the WikiText2~\cite{wikitext2} and C4~\cite{c4} datasets.

We also analyzed the runtime GPU memory usage and latency for both the \text{LLaMA LLM} and \text{BitNet b1.58}. The assessments utilized the FasterTransformer\footnote{\url{https://github.com/NVIDIA/FasterTransformer}} framework, which is highly optimized for GPU-based LLM inference latency. The 2-bit kernel from Ladder~\cite{ladder} was incorporated for \text{BitNet b1.58}. Our reported metric was the time taken per output token, as this is the primary factor affecting inference cost.

Table~\ref{tab:ppl} presents a summary of the perplexity scores and computational costs for both \text{BitNet b1.58} and \text{LLaMA LLM}. The data indicates that \text{BitNet b1.58} achieves comparable perplexity to the full precision \text{LLaMA LLM} starting at the 3B model size, while running 2.71 times faster and requiring 3.55 times less GPU memory. Notably, \text{BitNet b1.58} at 3.9B model size operates 2.4 times faster, uses 3.32 times less memory, and significantly outperforms the \text{LLaMA LLM} 3B model.

In Table~\ref{tab:zero-shot}, detailed zero-shot accuracy outcomes on various end tasks are shown. Our evaluation followed the protocol outlined by \emph{lm-evaluation-harness}\footnote{\url{https://github.com/EleutherAI/lm-evaluation-harness}}. Results reveal that the performance gap between \text{BitNet b1.58} and \text{LLaMA LLM} diminishes as model size increases. Crucially, \text{BitNet b1.58} matches the full precision baseline starting at the 3B scale. Consistent with the perplexity findings, the end-task results demonstrate that \text{BitNet b1.58} 3.9B surpasses \text{LLaMA LLM} 3B while demanding less memory and lower latency. This establishes \text{BitNet b1.58} as a Pareto improvement over current state-of-the-art LLMs.

\paragraph{Memory and Latency}

We expanded the model scale to 7B, 13B, and 70B parameters and assessed the associated costs. Figure~\ref{fig:latency-memory-energy} depicts the patterns of latency and memory usage, indicating that the speed-up improves as the model size increases. Specifically, the \text{BitNet b1.58} 70B model operates 4.1 times faster than the \text{LLaMA LLM} baseline. This improvement arises because the time required for \emph{nn.Linear} increases with model size. Memory usage exhibits a similar pattern, since the embedding layer remains in full precision and accounts for a smaller proportion of memory in larger models. Both latency and memory measurements were conducted using a 2-bit kernel, suggesting there remains potential for further optimization to decrease costs.

\paragraph{Energy}

We also evaluated the energy consumption of arithmetic operations for both \text{BitNet b1.58} and \text{LLaMA LLM}. Our focus was primarily on the matrix multiplication computations, as they represent the largest portion of LLMs' computational cost. Figure~\ref{fig:energy} breaks down the energy consumption components. Most operations in \text{BitNet b1.58} consist of INT8 additions, whereas \text{LLaMA LLM} performs both FP16 additions and multiplications. Based on the energy model described in~\cite{energycost, pokebnn}, \text{BitNet b1.58} reduces energy consumption for arithmetic operations in matrix multiplication by a factor of 71.4 on 7nm chips. We also provide end-to-end energy cost measurements for models processing 512 tokens. Our findings reveal that as model size grows, \text{BitNet b1.58} achieves progressively greater energy efficiency compared to the FP16 \text{LLaMA LLM} baseline. This efficiency gain is driven by the increasing share of \emph{nn.Linear} computations with model scaling, while the relative cost from other components diminishes for larger models.

\paragraph{Throughput}

We evaluate the throughput of \text{BitNet b1.58} compared to \text{LLaMA LLM} with 70 billion parameters running on two 80GB A100 GPUs, employing pipeline parallelism~\cite{gpipe} to accommodate the \text{LLaMA LLM} 70B model on the hardware. The batch size was progressively increased until reaching the GPU memory capacity, with a fixed sequence length of 512. As demonstrated in Table~\ref{tab:throughput}, \text{BitNet b1.58} 70B can handle a batch size up to 11 times larger than \text{LLaMA LLM}, yielding an 8.9-fold improvement in throughput.

\textbf{\text{BitNet b1.58} introduces a novel scaling law that relates model performance to inference cost}. For comparison, the equivalence across different model sizes between 1.58-bit and 16-bit precision can be established based on the outcomes shown in Figures~\ref{fig:latency-memory-energy} and \ref{fig:energy}.

\begin{itemize}
    \item The 13B \text{BitNet b1.58} model outperforms the 3B FP16 LLM in terms of latency, memory consumption, and energy efficiency.
    \item The 30B \text{BitNet b1.58} model is more efficient than the 7B FP16 LLM across latency, memory usage, and energy consumption metrics.
    \item The 70B \text{BitNet b1.58} model surpasses the 13B FP16 LLM regarding latency, memory footprint, and energy efficiency.
\end{itemize}

\paragraph{Training with 2T Tokens}

The volume of training tokens is a key parameter in LLM development. To assess the token scalability of \text{BitNet b1.58}, we trained a \text{BitNet b1.58} model on 2 trillion tokens following the data protocol of StableLM-3B~\cite{StableLM-3B-4E1T}, recognized as the leading open-source 3B model. Both models were evaluated using a benchmark suite composed of Winogrande~\cite{winoGrande}, PIQA~\cite{piqa}, SciQ~\cite{sciq}, LAMBADA~\cite{lambada}, and ARC-easy~\cite{arc}. Zero-shot accuracy results are presented in Table~\ref{tab:2t}. For tasks measured via accuracy and normalized accuracy, we report the average of both. The StableLM 3B results at 2T tokens were taken directly from its technical documentation. Our results reveal that \text{BitNet b1.58} attains superior performance across all evaluated tasks, indicating that 1.58-bit LLMs exhibit strong generalization capabilities as well.

\section{Discussion and Future Work}

\textbf{1-bit Mixture-of-Experts (MoE) LLMs}

Mixture-of-Experts (MoE) have demonstrated to be an efficient method for LLMs in terms of cost. Although it greatly reduces computational FLOPs, the considerable memory requirements and inter-chip communication overhead restrict its deployment and usability. These obstacles can be overcome by 1.58-bit LLMs. Primarily, the lowered memory demand decreases the number of devices needed to run MoE models. Additionally, it substantially cuts down the overhead involved in transferring activations across networks. Ultimately, if the entire model fits on a single chip, such overhead would be eliminated.

\textbf{Native Support of Long Sequence in LLMs}

With the rise of LLMs, supporting long sequences has become an essential capability. A significant hurdle for long sequence inference is the memory consumption caused by KV caches. \text{BitNet b1.58} marks a notable advancement towards native long sequence support by lowering activations from 16 bits to 8 bits, enabling the doubling of context length with the same resource availability. This can be further compressed losslessly to 4 bits or even less for 1.58-bit LLMs, which is left for future investigation.

\textbf{LLMs on Edge and Mobile}

Employing 1.58-bit LLMs can substantially enhance language model performance on edge and mobile platforms. These platforms often face constraints in memory and computational capacity, limiting LLM size and effectiveness. The lowered memory and power demands of 1.58-bit LLMs make their deployment feasible on such devices, unlocking numerous previously unattainable applications. This advancement can significantly boost the capabilities of edge and mobile devices and foster novel and innovative uses of LLMs. Furthermore, 1.58-bit LLMs are more compatible with CPU architectures, which dominate edge and mobile computing, allowing \text{BitNet b1.58} to run efficiently on these platforms and further enhance their performance.

\textbf{New Hardware for 1-bit LLMs}

Recent developments such as~Groq\footnote{\url{https://groq.com/}} have shown promising outcomes and strong potential in developing specialized hardware (e.g., LPUs) tailored for LLMs. Taking this further, we propose and encourage the design of new hardware and systems specifically optimized for 1-bit LLMs, leveraging the novel computational paradigm introduced by BitNet~\cite{bitnet}.

\end{document}